\subsection{Inhomogeneous Lorentz Group (Connected to the Identity)}\label{subsec:inhomogeneous-lorentz-group-(connected-to-the-identity)}
This original formulation of Axiom 5 (Lorentz Covariance) (\href{https://doi.org/10.1063/1.1704187}{Haag and Kastler}) uses the term ``inhomogeneous Lorentz group'' without specifying if this refers to all connected components of the inhomogeneous Lorentz group or simply the component of the inhomogeneous Lorentz group that is connected to the identity. This is likely a simple oversight.

The components of the inhomogeneous Lorentz group not connected to the identity are obtained from the component connected to the identity by reflections in time and/or space (See Section I.2.1 of  \href{https://doi.org/10.1007/978-3-642-61458-3}{Haag}), and it is an open question as to whether such reflections correspond to ``physical'' symmetries (Again see Section I.2.1 of  \href{https://doi.org/10.1007/978-3-642-61458-3}{Haag}).

That being said, it is a reasonable assumption to make that the phrase ``inhomogeneous Lorentz group'' used in Axiom 5 (Lorentz Covariance) should more properly be read as ``the component of the inhomogeneous Lorentz group connected to the identity''. We will adopt this posture.
